\section{An entropy inequality with one zero block}
\label{sec:block}

Our first objective is to separate the entropy produced by the block
pattern from the entropies within its three nonzero entries. The
rectangular formulation is needed because a channel output and its
Kraus index space can have different dimensions. The following is the
case of the triangular block entropy inequality of Tang, Zhu, Bai,
and Wang \cite[Theorem~4.2]{TangEtAl2026} in which all blocks have the
same rectangular shape, after permuting block rows. We give a direct
proof in the form needed for the thermal extension.

\begin{theorem}\label{thm:block}
Let $A,B,C\in\C^{d\times e}$, and put
\[
 M=\begin{pmatrix}A&B\\ C&0\end{pmatrix},\qquad
 x=\Tr(AA^*),\quad y=\Tr(BB^*),\quad z=\Tr(CC^*).
\]
If $x+y+z=1$, then
\begin{equation}\label{eq:block}
 \ent(MM^*)\ge
 x\ent(AA^*/x)+y\ent(BB^*/y)+z\ent(CC^*/z)+g(yz),
\end{equation}
where each term with zero weight is omitted.
\end{theorem}

The scalar matrix
$\bigl(\begin{smallmatrix}\sqrt x&\sqrt y\\\sqrt z&0\end{smallmatrix}\bigr)$
has squared singular values with sum $1$ and product $yz$.
The term $g(yz)$ is their entropy. To separate this contribution from
the entropies inside the three blocks, write the entropy deficit as
\[
 \Delta=\ent(AA^*)+\ent(BB^*)+\ent(CC^*)-\ent(MM^*).
\]
Since $\ent(AA^*)=x\ent(AA^*/x)-x\ln x$ for $x>0$, and similarly
for the other blocks, the theorem is equivalent to
\[
 \Delta\le H(x,y,z)-g(yz).
\]
We will express $\Delta$ as an integral of log-determinant differences.
Joint convexity then permits averaging over coordinate signs, and
homogeneous scalar concavity combines the resulting coordinates.
This integral and averaging route is also used in
\cite[Appendix~A.2--A.4]{TangEtAl2026}.

Block norm compression provides a related perspective. King
\cite[Theorem~1]{King2003} compared the Schatten norm of a positive
two-by-two block matrix with the norm of the scalar matrix of block
norms. Audenaert \cite[Conjecture~1]{Audenaert2008} studied the
corresponding comparison for arbitrary two-row block matrices and
proved several special cases. The entropy argument below allows
noncommuting blocks and does not require the general norm-compression
conjecture.

\subsection{From entropy to log determinants}

The direct sum of the three individual Gram matrices and the full Gram
matrix have equal trace. The following identity converts their entropy difference
into an integral, including when their dimensions or ranks differ.

\begin{lemma}\label{lem:integral}
If $U,V\ge0$ have equal trace, with possibly different sizes, then
\begin{equation}\label{eq:block-entropy-integral}
 \ent(U)-\ent(V)
 =\int_0^\infty
  \bigl[\ln\det(I+rU)-\ln\det(I+rV)\bigr]\frac{\dd r}{r^2}.
\end{equation}
The integral converges absolutely.
\end{lemma}

\begin{proof}
Let $f(r)$ denote the bracket. Equal trace gives $f(r)=O(r^2)$
at zero, while $f(r)=O(\ln r)$ at infinity. Hence the integral
converges absolutely and $f(r)/r$ vanishes at both endpoints.
Integration by parts reduces it to $\int_0^\infty f'(r)\dd r/r$.
If $a_i$ and $b_j$ are the positive eigenvalues of $U$ and $V$, then
\[
 \frac{f'(r)}r
 =\sum_i\frac{a_i}{r(1+ra_i)}
  -\sum_j\frac{b_j}{r(1+rb_j)}.
\]
A primitive of the term for $a>0$ is
$a\ln\bigl(r/(1+ra)\bigr)$. In the signed sum, its lower endpoint
limit is zero because the coefficients of $\ln r$ cancel by equal
trace. Its upper endpoint contribution is $-a\ln a$.
Taking their difference gives $\ent(U)-\ent(V)$. Zero eigenvalues
contribute nothing.
\end{proof}

\subsection{Joint convexity by a contraction}

The log-determinant function used in the averaging step is the following.
Its joint convexity follows by specialization and continuity from Hiai's
trace-function theorem \cite[preprint, Theorem~5.2]{Hiai2016}.
We give a direct proof: the contraction argument is the specialization
of \cite[Lemma~A.1]{TangEtAl2026} to a two-by-two block matrix.

\begin{lemma}\label{lem:logdet}
For each fixed $Z\in\C^{N\times N}$, the function
\begin{equation}\label{eq:block-G}
 G_Z(D,E)=\ln\det(D+ZE^{-1}Z^*)-\ln\det D,
 \qquad D,E>0,
\end{equation}
is jointly convex on pairs of positive definite Hermitian matrices.
\end{lemma}

\begin{proof}
Set
\[
 \mathcal D=\begin{pmatrix}D&0\\0&E\end{pmatrix},\qquad
 \mathcal C=\begin{pmatrix}D&Z\\-Z^*&E\end{pmatrix}.
\]
The block determinant identity gives
\[
 \det\mathcal C=\det E\det(D+ZE^{-1}Z^*)>0,
 \qquad G_Z(D,E)=\ln\det\mathcal C-\ln\det\mathcal D.
\]
Although $\mathcal C$ is not generally Hermitian, its determinant is
positive and its logarithm here is the real logarithm.
Along an affine line in $(D,E)$, let $\mathcal H$ be the derivative
of $\mathcal D$, which is also the derivative of $\mathcal C$.
The determinant derivative formula gives
\begin{equation}\label{eq:alt-second}
 G_Z''=\Tr(\mathcal D^{-1}\mathcal H\mathcal D^{-1}\mathcal H)
       -\Tr(\mathcal C^{-1}\mathcal H\mathcal C^{-1}\mathcal H).
\end{equation}
At the point under consideration put
\[
 L=\mathcal D^{-1/2}\mathcal H\mathcal D^{-1/2},\qquad
 K=\mathcal D^{-1/2}(\mathcal C-\mathcal D)\mathcal D^{-1/2}.
\]
Then $L=L^*$ and $K^*=-K$. Consequently
$\|(I+K)v\|_2^2=\|v\|_2^2+\|Kv\|_2^2$ for every vector $v$,
and $R=(I+K)^{-1}$ exists with operator norm at most one.
Cyclicity of trace rewrites \eqref{eq:alt-second} as
\[
 G_Z''=\|L\|_{\mathrm{HS}}^2-\Tr((RL)^2).
\]
The last trace is real because both log determinants along the line
are real. For any complex matrix $B$, entrywise Cauchy--Schwarz gives
$|\Tr(B^2)|=|\sum_{i,j}B_{ij}B_{ji}|\le\|B\|_{\mathrm{HS}}^2$.
Therefore
\[
 \Tr((RL)^2)\le|\Tr((RL)^2)|
 \le\|RL\|_{\mathrm{HS}}^2\le\|L\|_{\mathrm{HS}}^2.
\]
Every second directional derivative is thus nonnegative. The positive
definite domain is convex, so $G_Z$ is convex along each line segment,
as required.
\end{proof}

\subsection{The scalar deficit and its homogeneity}

We next use the same convexity to control the scalar entropy deficits
that sign averaging will produce. For $x,y,z\ge0$, define
\begin{align}
 m&=x+y+z,\qquad
 \lambda_\pm=\frac{m\pm\sqrt{m^2-4yz}}2,
 \label{eq:block-lambda}\\
 F(x,y,z)&=-x\ln x-y\ln y-z\ln z
           +\lambda_+\ln\lambda_++\lambda_-\ln\lambda_-.
 \label{eq:block-F}
\end{align}
Here $m^2\ge4yz$, and $\lambda_\pm$ are nonnegative with sum $m$
and product $yz$. Thus $F$ is continuous on the nonnegative octant,
with the usual convention at zero. On the probability simplex it is
exactly the deficit in our target: $F(x,y,z)=H(x,y,z)-g(yz)$.

\begin{lemma}\label{lem:scalar}
The function $F$ is positively homogeneous of degree one and concave
on the probability simplex. For every finite collection of nonnegative
triples,
\begin{equation}\label{eq:block-superadditive}
 \sum_i F(x_i,y_i,z_i)
 \le F\!\left(\sum_i x_i,\sum_i y_i,\sum_i z_i\right).
\end{equation}
\end{lemma}

\begin{proof}
The nonnegative entries in the two lists $(x,y,z)$ and
$(\lambda_+,\lambda_-)$ have the same sum. Applying
\Cref{lem:integral} to the corresponding diagonal matrices gives
\begin{align}
 F(x,y,z)
 &=\int_0^\infty
 \ln\frac{(1+rx)(1+ry)(1+rz)}
          {(1+r\lambda_+)(1+r\lambda_-)}\frac{\dd r}{r^2}\notag\\
 &=\int_0^\infty\left[\ln(1+rx)
 -\ln\!\left(1+\frac{rx}{(1+ry)(1+rz)}\right)\right]
          \frac{\dd r}{r^2}.
 \label{eq:scalar-deficit-integral}
\end{align}
The second equality uses
$(1+r\lambda_+)(1+r\lambda_-)=1+r(x+y+z)+r^2yz$;
the integral is absolutely convergent by the lemma.
For $x=1$ its bracket is
\[
 \ln(1+r)-G_{\sqrt r}(1+ry,1+rz).
\]
By \Cref{lem:logdet}, this is concave in $(y,z)$ for every $r>0$.
Integrating the concavity inequality shows that $F(1,\cdot,\cdot)$
is concave on the nonnegative quadrant.

Under scaling of $(x,y,z)$ by a positive constant, both
$\lambda_\pm$ scale by the same constant. Their sum is $m$, so the
logarithms of the scaling constant cancel in \eqref{eq:block-F}.
This proves homogeneity, including scale zero by continuity.
For a nonempty collection with all $x_i>0$, set $x=\sum_i x_i$.
Homogeneity and the preceding concavity now give
\begin{align*}
 \sum_iF(x_i,y_i,z_i)
 &=x\sum_i\frac{x_i}{x}F(1,y_i/x_i,z_i/x_i)\\
 &\le xF\!\left(1,\frac{\sum_i y_i}{x},
                     \frac{\sum_i z_i}{x}\right)
 =F\!\left(x,\sum_i y_i,\sum_i z_i\right).
\end{align*}
For a collection containing zero first coordinates, replace every
$x_i$ by $x_i+\varepsilon$, apply this inequality, and let
$\varepsilon\downarrow0$. Continuity of $F$ gives
\eqref{eq:block-superadditive}, even when every $x_i=0$.
The empty collection gives zero on both sides. Finally, applying
superadditivity to the two triples $\theta a$ and $(1-\theta)b$
and using homogeneity gives
$F(\theta a+(1-\theta)b)\ge\theta F(a)+(1-\theta)F(b)$
for nonnegative triples $a,b$ and $0\le\theta\le1$.
In particular, $F$ is concave on the probability simplex.
\end{proof}

\subsection{Sign averaging and the block estimate}

The two convexity consequences are now in place. The Schur complement
will put the matrix deficit into the form involving $G_Z$; averaging
over signs will replace its matrix arguments by diagonal entries;
and \Cref{lem:scalar} will combine the resulting scalar deficits.

\begin{proof}[Proof of \Cref{thm:block}]
First suppose that all three blocks are square of size $N$, and
temporarily allow arbitrary nonnegative masses $x,y,z$.
The matrices
\[
 U=AA^*\oplus BB^*\oplus CC^*,\qquad V=MM^*
\]
have equal trace. For $r>0$, put $D=I+rBB^*$ and $E=I+rC^*C$.
The Schur complement of $I+rCC^*$ in $I+rMM^*$ yields
\begin{equation}\label{eq:block-schur}
 \det(I+rMM^*)
 =\det(I+rCC^*)\det(D+rAE^{-1}A^*),
\end{equation}
where
$I-rC^*(I+rCC^*)^{-1}C=(I+rC^*C)^{-1}$.
Consequently the integrand bracket in \eqref{eq:block-entropy-integral}
is
\begin{equation}\label{eq:block-integrand}
 \ln\det(I+rAA^*)-G_{\sqrt r A}(D,E).
\end{equation}

Choose a singular-value decomposition $A=U_A\Sigma V_A^*$.
Replacing $(A,B,C)$ by $(\Sigma,U_A^*B,CV_A)$ amounts to
multiplying $M$ on the left and right by block-diagonal unitaries;
all entropies and masses are preserved. In these coordinates write
\[
 A=\diag(\sqrt{x_1},\ldots,\sqrt{x_N}),\qquad
 y_i=(BB^*)_{ii},\qquad z_i=(C^*C)_{ii}.
\]
These entries are nonnegative and satisfy
$\sum_i(x_i,y_i,z_i)=(x,y,z)$.
For each $\varepsilon\in\{-1,1\}^N$, let
$S_\varepsilon=\diag(\varepsilon_1,\ldots,\varepsilon_N)$.
Because $S_\varepsilon A S_\varepsilon=A$,
\[
 G_{\sqrt r A}(S_\varepsilon DS_\varepsilon,
                S_\varepsilon ES_\varepsilon)
 =G_{\sqrt r A}(D,E).
\]
Averaging simultaneously over all signs removes the off-diagonal
entries of both arguments. Joint convexity from \Cref{lem:logdet}
therefore gives
\[
 G_{\sqrt r A}(D,E)\ge
 G_{\sqrt r A}\bigl(\diag(1+ry_i),\diag(1+rz_i)\bigr).
\]
Because $G$ is subtracted in \eqref{eq:block-integrand}, that bracket
is bounded above by
\begin{align}
 &\sum_i\left[\ln(1+rx_i)
 -\ln\!\left(1+\frac{rx_i}{(1+ry_i)(1+rz_i)}\right)\right]\notag\\
 &\hspace{12mm}=
 \sum_i\ln\frac{(1+rx_i)(1+ry_i)(1+rz_i)}
                    {1+r(x_i+y_i+z_i)+r^2y_iz_i}.
 \label{eq:block-scalar-integrands}
\end{align}
By \eqref{eq:scalar-deficit-integral}, each summand integrates against
$\dd r/r^2$ to $F(x_i,y_i,z_i)$, and these finitely many integrals
converge absolutely. Integrating
the inequality and applying \Cref{lem:scalar} proves
\begin{equation}\label{eq:block-unnormalized}
 \ent(AA^*)+\ent(BB^*)+\ent(CC^*)-\ent(MM^*)
 \le\sum_iF(x_i,y_i,z_i)\le F(x,y,z).
\end{equation}
When $x+y+z=1$, the last expression is $H(x,y,z)-g(yz)$.
For $x>0$,
$\ent(AA^*)=x\ent(AA^*/x)-x\ln x$, and the same identity holds
for the other two blocks. Substitution gives \eqref{eq:block}.
If a mass vanishes, the corresponding block is zero, and its
contribution is zero. Thus the proof includes every boundary case;
all inverses used above are of positive definite matrices even
when $A,B,C$ are singular.

Finally, for $d\times e$ blocks let $N=\max(d,e)$ and choose the
coordinate isometries $J:\C^d\to\C^N$ and $K:\C^e\to\C^N$.
Replace each block $A$ by $JAK^*$, and likewise for $B,C$.
The resulting square block matrix is
$(J\oplus J)M(K\oplus K)^*$, so this operation only adds zero
singular values to $M$ and to each block. All masses and entropies
in \eqref{eq:block} are unchanged. The square case therefore proves
the rectangular case as well.
\end{proof}

\subsection{A scalar convexity fact for the channel application}

The scalar deficit estimate completes the zero-block theorem. The later
channel argument also needs a different property of the same entropy
function: its convexity as a function of the square root of the
determinant. We record it separately.

\begin{lemma}\label{lem:binary-spectrum}
The function $s\mapsto g(s^2)$ is nondecreasing and convex on $[0,1/2]$.
Consequently, $g$ is nondecreasing on $[0,1/4]$.
\end{lemma}

\begin{proof}
For $0<u<1/4$, put
\[
 t=\sqrt{1-4u},\qquad L=\frac{\atanh t}{t},\qquad
 \atanh t=\frac12\ln\frac{1+t}{1-t}.
\]
Differentiating the definition of $g$ gives
\begin{equation}\label{eq:g-derivatives}
 g'(u)=2L,\qquad
 g''(u)=\frac{4L-1/u}{t^2},\qquad
 ug''(u)+\frac12g'(u)=\frac{L-1}{t^2}\ge0.
\end{equation}
The last inequality follows from $\atanh t\ge t$, since
$(\atanh t)'=(1-t^2)^{-1}\ge1$ and both functions vanish at zero.
For $0<s<1/2$, these identities give
\[
 \frac{\mathrm d}{\mathrm ds}g(s^2)=2s g'(s^2)\ge0,\qquad
 \frac{\mathrm d^2}{\mathrm ds^2}g(s^2)
 =4\left(s^2g''(s^2)+\frac12g'(s^2)\right)\ge0.
\]
Continuity supplies both conclusions at the endpoints; no boundary
derivatives are required. Since $u\mapsto\sqrt u$ is nondecreasing,
$g(u)=g((\sqrt u)^2)$ is nondecreasing as well.
\end{proof}

The function $\kappa\mapsto g(\kappa^2/4)/\ln2$ is the
entropy--concurrence function of Wootters
\cite[Eq.~(8)]{Wootters1998}. Its monotonicity and convexity are the
scalar properties just proved. We use these properties, rather than
any additivity statement about entanglement of formation.
